\begin{postsub}[Post-submission addition: from the Euclidean case to general metric spaces]
The lemma above characterises the Euclidean conditional mean as the minimiser of a
\emph{weighted} least-squares criterion whose weights $s(z, x)$ depend on the predictor
only through the first two moments of $X$. The global Fréchet regression model of
\textcite{PetersenMüller2019} is obtained by keeping these weights and replacing the
Euclidean distance $d_E$ by a general metric $d$.
\begin{definition}[Global Fréchet regression]
    \label{def:global_frechet}
    Let $(\Omega, d)$ be a metric space, let $X \in \mathbb{R}^r$ be a predictor with mean
    $\mu$ and non-singular covariance matrix $\Sigma$, and define
    \[
        s(z, x) = 1 + (z - \mu)^\top \Sigma^{-1} (x - \mu).
    \]
    The \emph{global Fréchet regression function} is
    \begin{equation}
        \label{eq:global_frechet_pop}
        m_\oplus(x) = \argmin_{\omega \in \Omega}\,
        \mathbb{E}\!\left[ s(X, x)\, d^2(Y, \omega) \right].
    \end{equation}
    Given a sample $(X_i, Y_i)_{i=1}^n$, write $\bar{X}$ for the sample mean and
    $\hat{\Sigma}$ for the sample covariance of the $X_i$, and set
    \[
        s_{in}(x) = 1 + (X_i - \bar{X})^\top \hat{\Sigma}^{-1} (x - \bar{X}).
    \]
    The corresponding estimator is
    \begin{equation}
        \label{eq:global_frechet_emp}
        \hat{m}_\oplus(x) = \argmin_{\omega \in \Omega}\,
        \frac{1}{n} \sum_{i=1}^{n} s_{in}(x)\, d^2(Y_i, \omega).
    \end{equation}
\end{definition}
Setting $\Omega = \mathcal{G}$ and $d = d_W$ specialises this to our context. Because
$d_W = d_Q$ by Lemma~\ref{lemma:dqeqdw}, the criterion \eqref{eq:global_frechet_emp}
becomes a weighted average of squared $\mathcal{L}^2$ distances between quantile
functions,
\begin{equation}
    \label{eq:frechet_qf_objective}
    \hat{m}_\oplus(x) = \argmin_{Q}\,
    \frac{1}{n} \sum_{i=1}^{n} s_{in}(x) \int_0^1 \bigl(Q_i(u) - Q(u)\bigr)^2 \, du,
\end{equation}
whose \emph{unconstrained} minimiser is the weighted cross-sectional mean of the sample
quantile functions,
\[
    \tilde{Q}(x)(\cdot) = \frac{1}{n} \sum_{i=1}^{n} s_{in}(x)\, Q_i(\cdot).
\]
When $\tilde{Q}(x)$ is non-decreasing it is already a valid quantile function and hence the
solution; otherwise $\hat{m}_\oplus(x)$ is its $\mathcal{L}^2$-projection onto the cone of
non-decreasing functions, which is precisely the quadratic program in
Algorithm~\ref{alg:quadprog}.
\end{postsub}
